% Dutch edition additions. Loaded after layout/hp-book, before \begin{document}.
\usepackage{polyglossia}
\setmainlanguage{dutch}
\setotherlanguage{english}

% Inline gloss for the A2-B1 edition: \gl{woord}{english} -> woord (english)
\newcommand{\gl}[2]{#1 (\textit{#2})}

% fmtcount (used by hp-book.sty for "CHAPTER ONE") has no Dutch, so the chapter
% number word is supplied by nl-assemble as \nlchapterword before each \chapter.
\providecommand{\nlchapterword}{\thechapter}
\makeatletter
\renewcommand*{\printchapternum}{%
  \chapnumfont \hyperref[contents]{\textls[30]{\nlchapterword}}}
\renewcommand*{\printchaptername}{%
  \centering \chapnamefont \hyperref[contents]{\MakeUppercase{\textls[30]{Hoofdstuk}}}}
\makeatother

\hyphenation{Zwein-stein Per-ka-men-tus An-der-ling Grif-foen-dor Zwad-de-rich
Ra-ven-klauw Huf-fel-puf Dreu-zel Dreu-zels Wegis-weg Goud-grijp Mal-fi-dus
Her-me-lien Grif-fel Lub-ber-mans Krin-kel Pot-ter-Evans-Ver-res Ver-res-Evans}
